% GENERATED FILE. Edit canonical lessons, then run npm run generate:books.
% canonical-source-hash: fnv1a64:c8b6d5102d7e8846
% canonical-lessons: LA-C332-heres, LA-C332-hereditas, LA-C332-testamentum, LA-C332-obsigno, LA-C332-relinquo

\chapter{Testāmentum}
\label{ch:la-testamentum}
\hlchaptermodality{332}

\begin{chapteropening}
\textbf{By the end of this chapter:} \emph{I can talk about a will and inheritance: an heir, an inheritance, a will, to seal and to leave behind.}

Complete the last lesson of chapter 332: I can talk about a will and inheritance: an heir, an inheritance, a will, to seal and to leave behind.
\end{chapteropening}

\section[hērēs, hērēdis]{hērēs, hērēdis — an heir}

\label{lesson:LA-C332-heres}



\begin{quote}
Before the new one: say the Latin for a memorial, then the Latin for offerings to the dead.
\end{quote}

\subsection*{You'll want to know: hērēs, hērēdis}
\textbf{hērēs, hērēdis} — \textquotedblleft{}an heir\textquotedblright{}.

\begin{quote}
Fīlius ūnus hērēs patris erat.
\end{quote}

\begin{quote}
\emph{The only son was his father's heir.}
\end{quote}

Masculine or feminine, 3rd declension. \emph{Hērēdem facere}: to name someone heir in a will. English \emph{heir}.

\begin{cousinweb}
\textbf{hērēs} is inherited; EtymDB traces it to a Proto-Indo-European root meaning \textquotedblleft{}to leave behind\textquotedblright{}, while Lewis \& Short compare a Sanskrit root meaning \textquotedblleft{}to take\textquotedblright{}, so the details are not settled. English \emph{heir} comes from it through French.
\end{cousinweb}

\begin{grammarlens}[title={What you've built}]
One more word for telling, explaining and describing at B1.
\end{grammarlens}

\subsection*{Your turn}
\begin{itemize}
\raggedright
  \item \emph{Say it:} \emph{hērēs, hērēdis}
  \item \emph{Say it:} \emph{hērēs, hērēdis}, once more
  \item \emph{Say it:} a sentence of your own with \emph{hērēs, hērēdis}
  \item \emph{Recall:} say the Latin for a memorial, then the Latin for offerings to the dead, then say \emph{hērēs, hērēdis} again
\end{itemize}

\begin{tcolorbox}[breakable,colback=teal!4,colframe=teal!35!black,title={Before you move on}]
What does hērēs, hērēdis mean? (\textquotedblleft{}An heir\textquotedblright{}.) Say it once more.
\end{tcolorbox}
\section[hērēditās, hērēditātis]{hērēditās, hērēditātis — an inheritance}

\label{lesson:LA-C332-hereditas}



\begin{quote}
Before the new one: say the Latin for offerings to the dead, then the Latin for an heir.
\end{quote}

\subsection*{You'll want to know: hērēditās, hērēditātis}
\textbf{hērēditās, hērēditātis} — \textquotedblleft{}an inheritance\textquotedblright{}.

\begin{quote}
Magnam hērēditātem ab avunculō accēpit.
\end{quote}

\begin{quote}
\emph{He received a large inheritance from his uncle.}
\end{quote}

Feminine, 3rd declension, from \emph{hērēs}. English \emph{heredity} and \emph{heritage} belong to its family.

\begin{cousinweb}
\textbf{hērēditās} is built on \emph{hērēs}, \textquotedblleft{}heir\textquotedblright{} (EtymDB), with the ending \emph{-tās}: what comes to an heir. English \emph{heredity} comes from it.
\end{cousinweb}

\begin{grammarlens}[title={What you've built}]
One more word for telling, explaining and describing at B1.
\end{grammarlens}

\subsection*{Your turn}
\begin{itemize}
\raggedright
  \item \emph{Say it:} \emph{hērēditās, hērēditātis}
  \item \emph{Say it:} \emph{hērēditās, hērēditātis}, once more
  \item \emph{Say it:} a sentence of your own with \emph{hērēditās, hērēditātis}
  \item \emph{Recall:} say the Latin for offerings to the dead, then the Latin for an heir, then say \emph{hērēditās, hērēditātis} again
\end{itemize}

\begin{tcolorbox}[breakable,colback=teal!4,colframe=teal!35!black,title={Before you move on}]
What does hērēditās, hērēditātis mean? (\textquotedblleft{}An inheritance\textquotedblright{}.) Say it once more.
\end{tcolorbox}
\section[testāmentum, testāmentī]{testāmentum, testāmentī — a will, a testament}

\label{lesson:LA-C332-testamentum}



\begin{quote}
Before the new one: say the Latin for an heir, then the Latin for an inheritance.
\end{quote}

\subsection*{You'll want to know: testāmentum, testāmentī}
\textbf{testāmentum, testāmentī} — \textquotedblleft{}a will, a testament\textquotedblright{}.

\begin{quote}
Senex testāmentum cōram testibus scrīpsit.
\end{quote}

\begin{quote}
\emph{The old man wrote his will in front of witnesses.}
\end{quote}

Neuter, from a verb meaning 'to call to witness', and so from \emph{testis} (taught). Seven witnesses sealed a Roman will.

\begin{cousinweb}
\textbf{testāmentum} is built on \emph{testor}, \textquotedblleft{}to call to witness\textquotedblright{}, from \emph{testis}, \textquotedblleft{}witness\textquotedblright{} (EtymDB): a declaration made before witnesses. English \emph{testament} comes from it.
\end{cousinweb}

\begin{grammarlens}[title={What you've built}]
One more word for telling, explaining and describing at B1.
\end{grammarlens}

\subsection*{Your turn}
\begin{itemize}
\raggedright
  \item \emph{Say it:} \emph{testāmentum, testāmentī}
  \item \emph{Say it:} \emph{testāmentum, testāmentī}, once more
  \item \emph{Say it:} a sentence of your own with \emph{testāmentum, testāmentī}
  \item \emph{Recall:} say the Latin for an heir, then the Latin for an inheritance, then say \emph{testāmentum, testāmentī} again
\end{itemize}

\begin{tcolorbox}[breakable,colback=teal!4,colframe=teal!35!black,title={Before you move on}]
What does testāmentum, testāmentī mean? (\textquotedblleft{}A will, a testament\textquotedblright{}.) Say it once more.
\end{tcolorbox}
\section[obsignō, obsignāre]{obsignō, obsignāre — to seal, to seal up}

\label{lesson:LA-C332-obsigno}



\begin{quote}
Before the new one: say the Latin for an inheritance, then the Latin for a will.
\end{quote}

\subsection*{You'll want to know: obsignō, obsignāre}
\textbf{obsignō, obsignāre} — \textquotedblleft{}to seal, to seal up\textquotedblright{}.

\begin{quote}
Testēs testāmentum ānulīs suīs obsignāvērunt.
\end{quote}

\begin{quote}
\emph{The witnesses sealed the will with their rings.}
\end{quote}

\emph{obsignō, obsignāre, obsignāvī, obsignātum}: \emph{ob} + \emph{signō} (taught). A seal of wax pressed with a signet ring.

\begin{cousinweb}
\textbf{obsignō} is built on \emph{signō}, \textquotedblleft{}to mark, to seal\textquotedblright{}, with \emph{ob-} (EtymDB): to put a seal over.
\end{cousinweb}

\begin{grammarlens}[title={What you've built}]
One more word for telling, explaining and describing at B1.
\end{grammarlens}

\subsection*{Your turn}
\begin{itemize}
\raggedright
  \item \emph{Say it:} \emph{obsignō, obsignāre}
  \item \emph{Say it:} \emph{obsignō, obsignāre}, once more
  \item \emph{Say it:} a sentence of your own with \emph{obsignō, obsignāre}
  \item \emph{Recall:} say the Latin for an inheritance, then the Latin for a will, then say \emph{obsignō, obsignāre} again
\end{itemize}

\begin{tcolorbox}[breakable,colback=teal!4,colframe=teal!35!black,title={Before you move on}]
What does obsignō, obsignāre mean? (\textquotedblleft{}To seal, to seal up\textquotedblright{}.) Say it once more.
\end{tcolorbox}
\section[relinquō, relinquere]{relinquō, relinquere — to leave behind; to leave (in a will)}

\label{lesson:LA-C332-relinquo}



\begin{quote}
Before the new one: say the Latin for a will, then the Latin for to seal.
\end{quote}

\subsection*{You'll want to know: relinquō, relinquere}
\textbf{relinquō, relinquere} — \textquotedblleft{}to leave behind; to leave (in a will)\textquotedblright{}.

\begin{quote}
Pater fīliīs domum et agrōs relīquit.
\end{quote}

\begin{quote}
\emph{The father left the house and the fields to his sons.}
\end{quote}

\emph{relinquō, relinquere, relīquī, relictum}. English \emph{relinquish} and \emph{relic} belong to its family.

\begin{cousinweb}
\textbf{relinquō} is built on \emph{linquō}, \textquotedblleft{}to leave\textquotedblright{}, with \emph{re-}, \textquotedblleft{}back, behind\textquotedblright{} (EtymDB). English \emph{relinquish} comes from it through French, and \emph{relic} from \emph{reliquiae}, \textquotedblleft{}remains\textquotedblright{}, of the same family.
\end{cousinweb}

\begin{grammarlens}[title={What you've built}]
One more word for telling, explaining and describing at B1.
\end{grammarlens}

\subsection*{Your turn}
\begin{itemize}
\raggedright
  \item \emph{Say it:} \emph{relinquō, relinquere}
  \item \emph{Say it:} \emph{relinquō, relinquere}, once more
  \item \emph{Say it:} a sentence of your own with \emph{relinquō, relinquere}
  \item \emph{Recall:} say the Latin for a will, then the Latin for to seal, then say \emph{relinquō, relinquere} again
\end{itemize}

\begin{tcolorbox}[breakable,colback=teal!4,colframe=teal!35!black,title={Before you move on}]
What does relinquō, relinquere mean? (\textquotedblleft{}To leave behind; to leave (in a will)\textquotedblright{}.) Say it once more.
\end{tcolorbox}
